% SPDX-License-Identifier: Apache-2.0
\section{\code{txhdl\_macros}}
\label{sec:r-macros}

The procedural macros, written against \code{proc\_macro} alone so the
crate has no dependencies. \code{\#[derive(Transaction)]} and
\code{\#[derive(Bus)]} emit one empty \code{impl} each.
\code{\#[derive(Value)]} gives a struct of values, or a fieldless
enum, its width and its bits for a trace, and \code{\#[derive(Trace)]}
registers every field of a unit under its field name; both read the
item's generic parameters so a generic unit can derive them.
\code{\#[pipeline(op = latency, ..)]} is the first of the lowering: on
an \code{async fn} of awaited operator calls it keeps the function and
emits a \code{NAME\_verilog()} beside it, each operator a stage of the
stated latency and each value delayed to meet the operator that uses
it. Its grammar is the one \code{mac} needs and no wider, and it
refuses the rest with its own message. \code{\#[lower]} on an
\code{impl Unit for Unit} does the same for a unit with state,
filling the header's \code{<In, Out>} in from \code{run}'s signature,
so the ports are named once: it reads \code{run} as one loop, or
several under
\code{join2}, each a process that waits for the rising or the falling
edge of its clock, reads registers and ports into names, predicates
drives with
\code{with!}, \code{case!} and \code{if}, chooses values
with
\code{select!},
and drives registers, words of memories and output ports, and writes
\code{Unit::lowered(name)}, \code{verilog} and \code{vhdl} beside
the impl. Ports come from \code{run}'s signature, a side of several
written as a tuple or as a struct whose fields are the ports;
registers, memories and widths come from the
struct through the \code{Trace} derive, which also derives
\code{Fields} and \code{Port}; and the configuration's constants are
evaluated when \code{lowered} runs, so a generic unit lowers once per
build. A \code{let} that computes is a wire named for it, a read of a
port or register an alias. A call of a function under \code{\#[lower]}
is inlined in the same terms: each argument and each \code{let} of the
function is bound to the expression it stands for, and one the body
reads more than once, and that is worth more than a name, is given a
wire of its own instead, called after the function and the name it had
there and numbered. That keeps an expression written once whatever the
body does with it, so calls that nest cost a wire each rather than a
power of the reads, which is issue 126.
A helper that is not in the unit's file is reached another way, since
the file is all the macro reads. \code{\#[lower]} on a function adds,
beside it, a type of the same name in the other namespace, whose
\code{lowered} builds the function's expression when a unit is
lowered; a call the file does not explain, \code{f(x)} or
\code{m::f::<K>(x)}, lowers to \code{f::lowered(x)} or
\code{m::f::lowered::<K>(x)}, which Rust resolves wherever it resolves
the function, through a path or a \code{use}. Inside it a value read
twice is bound by \code{inline\_bind} to a wire of the unit, which
the unit's \code{lowered} collects between \code{inlined\_begin} and
\code{inlined\_end}, a frame per unit so that a child lowered inside
it keeps its own. A parameter typed as a struct is known by its type
while the body is read, so \code{p.f} is a field of it, whatever
computed the value; and a struct literal passed for it, or given back,
is the concatenation of its fields (issue 504). It reads \code{case!} and \code{select!},
and a \code{match} that yields a value, as chains on the value, with
an enum variant, an integer, a range \code{lo..=hi} or \code{lo..hi}
or \code{\_} as a pattern and a guard joined by \code{\&\&}; a
\code{match} splits its own arms, since rustfmt ends an arm whose
value is a block with no comma. An \code{if} that yields a value is the
same chain on its conditions, and is refused without an \code{else};
a branch or an arm is one expression. \code{(lo..=hi).contains(\&x)}
is the two compares a range pattern is. \code{*} and \code{\%} lower as
the netlist's, \code{\%} as \code{rem} in VHDL; unary \code{-} as
\code{\~{}x + 1} at \code{x}'s width, a negative number being refused
for having none; and \code{x as uN} as \code{Expr::cast}, which the
emitters resolve once the width is known, to \code{x} when it is no
wider and to its low bits when it is. It used to drop the cast, which
every one of the 132 casts in the tree survived only by widening
(issue 496); \code{slice},
\code{concat}, \code{sext} and \code{zext} by their turbofish, every
width it needs written, since it sees tokens and not types, and only
\code{concat}'s low operand left to Rust as \code{\_}; a
struct literal as the concatenation of its fields; an \code{if} chain
as arms under their conditions, each arm's statements lowered in
turn, a \code{send} in an arm hoisted with the arm's path as
\code{valid}, and a \code{set} of an output under it refused, since
an output is a wire; \code{self.m.at(a)}
on the left of a drive as a word of a memory and \code{self.m.read(a)}
as an index; and channels as a valid/ready handshake: \code{until} or
a channel's \code{wait} is the guard on every register drive after
it, a receive asserts \code{ready} under the guard, a \code{take}
binds \code{valid} and \code{data} to its two names and asserts
\code{ready} the same way, a send drives
\code{data} and asserts \code{valid} under it. What it does not read
it refuses with a message that names the construct.
A field of a port's value, \code{p.tag} where \code{p} is what a
channel port transports, is read as the slice of the port's data the
field occupies: the \code{Value} derive lays a struct's fields out
first field highest and says so in a \code{layout}, and the netlist
finds the bits when \code{lowered} runs, so a generic field's width
is no obstacle.
A side that is a struct, \code{inp: S<..>} or
\code{S \{ a, b \}: S<..>}, is read from the file the unit is in, as
the inlined functions are, since a proc macro sees only the item it
is on. \code{find\_port\_structs} collects every struct of the file,
with its generic parameters and its fields. It also collects every
role of an \code{interface!} in the file: a member of \code{Signal}
becomes an \code{Out} or an \code{In}, and a member of \code{Chan} a
\code{Tx} or an \code{Rx}, by the role's direction. The side's
arguments are substituted for the struct's parameters in each field's
type, and each field is a port named for the field, in declaration
order. Before the body is read, \code{port\_fields} rewrites every
\code{inp.a} of such a side to the port \code{a}, so the rest of the
lowering sees the names it sees for a tuple. A type that is neither a
port nor a struct of the file is taken to implement \code{Ports},
which \code{\#[derive(Ports)]} writes, so a struct declared in any
file of any crate can be a side (issue 483). Its fields cannot be
seen here, so \code{port\_fields} rewrites each \code{bus.f} the body
reads to the port \code{bus\_f}, named for the side as well as the
field so that two sides of one type stay apart. The ports come from
\code{bundle\_ports} when \code{lowered} runs, each field's kind,
width and clock from its type through \code{PortEnd}, and a field of
a port's value is found by \code{bundle\_field} from the same
layout. A struct of ports may hold another: \code{PortField} is what
a field is, one port for an end and every port for a struct of ports,
each \code{field\_sub}, so the netlist sees them flattened, and
\code{port\_fields} reads \code{bus.pins.awid} as the port
\code{bus\_pins\_awid}, every name of the chain up to a method call
being the port's. A struct of the unit's own file may hold one the
same way, and a side passed whole joins each nested struct whole
(issue 498). One rule names the ports of either kind of side: written
\code{name: S} its ports are \code{name\_field}, and taken apart,
\code{S \{ a, b \}: S}, which only a struct of the file allows, they
are \code{a} and \code{b} (issue 580). \code{instance} joins a parent's struct literal to such a
child by the one port that ends in the field's name. A unit of units
may take such a side and pass its fields, \code{bus.req}, to its
children, which join as the parent's ports; their kinds cannot be
seen here, so \code{Lowered::checked} checks when \code{lowered}
runs that each channel port of a unit of units is joined to exactly
one child. Passing the whole side to a child joins its ports to the
child's in the order the struct declares them, and since the macro
cannot see them the joins are made when \code{lowered} runs, by
\code{bundle\_args} and \code{instance\_of} (issue 498). A port name
that two sides both give is refused.
An array of ports, \code{[Rx<T>; N]}, is such a side: \code{Ports} is
implemented for arrays, their fields the indices, so its ports are
\code{ins\_0} to \code{ins\_\{N-1\}} and are found when \code{lowered}
runs. A \code{for i in lo..hi} at the top of the loop is unrolled
there too, since \code{N} is a const parameter the macro cannot see
(issue 500). A process's body is written as a block that pushes its
statements, so a loop is a Rust \code{for} among the pushes;
\code{port\_fields} writes \code{ins[i]} as a name holding \code{i}
between two marks, and \code{dyn\_names} turns every quoted name
holding one into a \code{format!} of it in the text the macro writes.
Inside a loop, and for a \code{let mut}, a \code{let} is a Rust
variable of the expression, and each value computed is a wire of its
own, pushed onto \code{\_\_dynw} and named \code{x\_lN} by a counter,
so a value an assignment forwards from one turn to the next is written
once a turn. The process bodies are built before the unit, so those
wires join the others. A \code{for} under an \code{if} is refused.
So is a \code{let} after a loop that reads what the loop computed,
since a plain wire is defined outside the process and those variables
are not. A \code{send} under an \code{if} inside a \code{for} stays in
the loop, one drive a turn, where elsewhere it is hoisted to the end
of the process. An array may be one port among others in a tuple,
and a unit of units passes an array port its ends as
\code{[a, b, c]}, joined in order to \code{x\_0} to \code{x\_2}. An
index is a number where \code{lowered} runs, so a constant may use
it: \code{U::<3>::from(i)} is a literal on each turn.
A register of a \code{Regs} field, \code{self.prio[i]}, is the
register \code{prio\_i}, read or driven, and in \code{with!} and
\code{when!} it is \code{prio[i]: v}; the index is a number or a
loop's variable, since a register chosen by a signal is a
multiplexer to write out, and that is refused.
A name VHDL or SystemVerilog reserves, \code{next}, \code{out} and
\code{inside} among them and VHDL's without regard to case, is
escaped rather than refused: the netlist writes it with \code{\_rw}
after it, the same in both targets, whichever one reserves it (issue
497). A field is escaped by the \code{Trace} derive, which gives the
netlist and the trace the same name for it, through the table
\code{\#[rename]} already fills. Everything else is escaped by
\code{Lowered::escaped}, which the emitters and the ports file call
before they write: the module's own name, a string given to
\code{lowered} that the macro never sees; a port of one wire; a wire
or a net; an instance. A port escaped is recorded in the ports file
with its own name as its trace scope, so the testbench reads the
run's \code{next} for the netlist's \code{next\_rw}, and
\code{trace\_as} still takes a port's own name. A channel is left
alone, since its nets are its name and \code{\_data}, \code{\_valid}
or \code{\_ready}, and so is a foreign module, whose names are its
vendor's. An escaped name another name already has is refused,
naming both, which is the one case left. The two word lists are one
file, \code{lib/src/reserved.rs}, below: the runtime compiles it as
a module of its own, and the macro crate, which cannot depend on the
runtime since the runtime depends on it, includes the same file by
path, as does \code{//tools/fst2tb}, which escapes the entity it is
asked for the same way. A computed \code{let} or a
port whose netlist name is also a field's would be declared twice.
\code{\#[lower]} does not see the struct, so for each such name it
writes a constant, placed at the name, that tests the name against the
\code{NAMES} the derive gives \code{Fields} and fails to compile if it
is there; \code{lowered} uses every one. For a generic unit that
constant is evaluated only for a type that is lowered. A computed
\code{let} with a port's name is still renamed, with \code{\_w}, and a
name bound twice is numbered.
\code{station!(Name, N)} writes a reservation
station of \code{N} inputs as text, a line struct, the station's
struct and its lowered \code{run}, every wire of the step named for
what it is, and \code{TXHDL\_MACRO\_DUMP} writes the text out for
reading; the attribute in its output expands as any does.
A \code{run} with no loop is a unit of units, and \code{\#[lower]}
reads it as the hierarchy example writes it for the simulation: a
\code{let (tx, rx) = chan::<T, C>()} or \code{let (out, inp) =
signal::<T, C>()} is a net between the children, named by the two
ends' common prefix, and every \code{self.child.run(inputs,
outputs)}, alone or under \code{join2}, is an instance of the
child's own lowering, its ports joined in order to the nets and to
the parent's ports. A wire's reading end or an input port may be read
by several children: \code{let b = a.clone();} is another name for the
same net, which is how a reset reaches the core, the timer and the
serial port at once; a channel end is not cloned this way, since a
channel has one receiver. A side of the parent's that is a struct, passed
whole, is its fields in order, and a struct literal
\code{S \{ b: x, a \}} joins each field to the child's port of that
name, whatever order the literal writes them in; a channel end or
channel port joined twice or never is refused by name, as is anything else in the body.
A field of the literal that is itself a struct of ports takes a
literal of its own, \code{pins: P \{ awid, .. \}}, which joins the
child's \code{pins\_awid}, or the parent's own nested struct whole,
\code{pins: jtag}, which joins the child's \code{pins\_} ports to the
parent's \code{jtag\_} ones in order through \code{bundle\_args\_named}
when \code{lowered} runs (issue 579).
A unit of units may also join an array of children,
\code{join\_all(self.f.iter\_mut().enumerate().map(|(i, c)| c.run(..)))}:
the lowering writes a loop that pushes an instance per child,
\code{f\_i}, when \code{lowered} runs, and each \code{x.take(e)} among
the arguments joins the net or port \code{x\_e}, \code{e} being any
expression of \code{i}. The nets of \code{chans} and \code{signals} are
declared by a loop too, and \code{Ends::from} over an array port names
its elements; the ends are bound under names of their own, since a
side's name followed by a field is read as one of its ports (issue
635).
A child of the array may be passed \code{tie(v)} as well, where \code{v}
may be an expression of \code{i}, and each child then has wires of
its own, \code{f\_i\_tieK}; \code{Mesh} ties each node's column and row
that way (issue 635).
An input passed \code{tie(v)} is held at the constant \code{v}: in
the run it is a wire \code{comp::tie} sets once, and in the netlist
a wire of the parent, \code{FIELD\_tieK}, driven by \code{lit(v)},
so \code{v} is evaluated when \code{lowered} runs and is a constant.
\code{Lowered::checked} refuses such a wire joined to anything but a
child's input, which it would drive twice (issue 498). The
parent's \code{lowered} reaches each child through a value of
itself, so a unit of units is \code{Default}, as every unit is, and
each child's module is named for the parent and the field. A net
made by \code{chan} is the runtime's channel in the netlist, a
buffer of two with the receiver's \code{valid} and \code{data} and
the sender's \code{ready} registered, one module \code{txhdl\_chan}
put into the file once; a net made by \code{signal} is a wire. A
bundle is made whole by \code{let (host, per) = link::<B>()}, where
\code{B} is one side's struct of ports and implements \code{Link},
whose \code{Host} is the other side's: its nets are listed by
\code{link\_nets} when \code{lowered} runs, one per port of
\code{B}, each named for the net and the field, and a side passed
whole joins them in order through \code{link\_args}, or one at a time
as \code{host.aw}. \code{Lowered::checked} then counts every channel
between children, which must join exactly two ports (issue 498).
A child need not be lowered. A unit that implements \code{Lower} by
hand with \code{netlist::foreign} is a module the netlist does not
write: its lowering names the module, its ports by their own names in
the order the unit's \code{run} takes them, its parameters, and which
clock drives each of its clock pins, and the parent joins it as it
joins any child. The netlist then has an instance of that module,
with its parameters, where the child would be, and emits no body for
it: a Verilog instance with \code{\#(..)}, and in VHDL a component
declared once per module and instantiated with a generic map, since a
Verilog module is reached from VHDL no other way. The component's
vectors are \code{std\_logic\_vector} where the netlist's are
\code{unsigned}, so a vector crosses with a conversion, on the actual
for an input and on the formal for an output. A port of kind
\code{Pad} is a pin driven from both sides: the parent takes it among
its own ports and passes it to the foreign child untouched, and both
netlists declare it \code{inout}. The simulation transfers no signal state
across a pad; the foreign unit's model in Rust stands in for whatever is on
the other side of the pins.
\code{interface!} takes an interface and any number of roles, checks
that every role names every member once and that no member is driven
twice, and emits a struct per role plus a constructor that splits each
member. \code{with!(self <= \{ .. \})} is the drives of one struct,
its name written once: an entry \code{field: value} is a drive,
\code{c ? field: value} a drive under a condition, \code{c ? \{ ..
\} else \{ .. \}} a group under one with the entries after
\code{else} under its failure, and a memory word \code{m.at(i)} a
field like any other; entries apply in order, the last drive of a
field winning, and each lowers as an \code{if} in the clocked block;
an entry that is not \code{field: value} is refused with the macro's
own message. \code{when!(c => self \{ .. \} else \{ .. \})} is one
group of it with the condition written first, the same entries and
the same \code{if}; an \code{if} whose condition is a constant of
the build is folded, its arm kept or dropped and the test gone. \code{case!} is a predicated drive per arm that reads
as a conditional: each statement inside an arm is \code{register <=
value}, the arrow pointing into the register, and a statement that
is not is refused with the macro's own message. Both are procedural
because a \code{macro\_rules!} \code{expr} fragment may be followed
only by \code{=>}, \code{,} or \code{;}, so \code{lhs <= rhs} cannot be
matched declaratively at all.

\lstinputlisting[caption={\code{lib/macros/lib.rs}.},label={lst:r-macros}]{lib/macros/lib.rs}

\lstinputlisting[caption={\code{lib/src/reserved.rs}: the words the two targets reserve, shared by the macros and the runtime.},label={lst:r-reserved}]{lib/src/reserved.rs}
